% GENERATED FILE. Edit canonical lessons, then run npm run generate:books.
% canonical-source-hash: fnv1a64:a69a9383d818ca91
% canonical-lessons: PT-C479-magoar, PT-C479-magoa, PT-C479-rancor, PT-C479-perdoar, PT-C479-perdao

\chapter{Mágoas e perdão}
\label{ch:pt-magoas-e-perdao}
\hlchaptermodality{479}

\begin{chapteropening}
\textbf{By the end of this chapter:} \emph{I can talk about hurt, grudges and forgiveness.}

Complete the last lesson of chapter 479: I can talk about hurt, grudges and forgiveness.
\end{chapteropening}

\section[magoar]{magoar — to hurt (someone's feelings or body)}

\label{lesson:PT-C479-magoar}



\begin{quote}
Before the new one: say the Portuguese for lonely, then the Portuguese for grief.
\end{quote}

\subsection*{You'll want to know: magoar}
\textbf{magoar} — \textquotedblleft{}to hurt (someone's feelings or body)\textquotedblright{}.

\begin{quote}
As palavras dele magoaram-me muito.
\end{quote}

\begin{quote}
\emph{His words hurt me a lot.}
\end{quote}

Also physical: \emph{magoei o pé} — I hurt my foot. \emph{Magoar-se} — to get hurt.

\begin{cousinweb}
\emph{magoar} is usually traced to Latin \emph{maculāre} \textquotedblleft{}to spot, to stain\textquotedblright{} (Lewis \& Short), built on \emph{macula} \textquotedblleft{}a spot, a stain\textquotedblright{}: a hurt leaves a mark.
\end{cousinweb}

\begin{grammarlens}[title={What you've built}]
One more word for telling, explaining and describing at B1.
\end{grammarlens}

\subsection*{Your turn}
\begin{itemize}
\raggedright
  \item \emph{Say it:} \emph{magoar}
  \item \emph{Say it:} \emph{magoar}, once more
  \item \emph{Say it:} a sentence of your own with \emph{magoar}
  \item \emph{Recall:} say the Portuguese for lonely, then the Portuguese for grief, then say \emph{magoar} again
\end{itemize}

\begin{tcolorbox}[breakable,colback=teal!4,colframe=teal!35!black,title={Before you move on}]
What does magoar mean? (\textquotedblleft{}To hurt (someone's feelings or body)\textquotedblright{}.) Say it once more.
\end{tcolorbox}
\section[a mágoa]{a mágoa — hurt, sorrow}

\label{lesson:PT-C479-magoa}



\begin{quote}
Before the new one: say the Portuguese for grief, then the Portuguese for to hurt.
\end{quote}

\subsection*{You'll want to know: a mágoa}
\textbf{a mágoa} — \textquotedblleft{}hurt, sorrow\textquotedblright{}.

\begin{quote}
Ainda guarda mágoa da discussão com o irmão.
\end{quote}

\begin{quote}
\emph{She still feels hurt about the argument with her brother.}
\end{quote}

\emph{Guardar mágoa} — to keep feeling hurt.

\begin{cousinweb}
\emph{mágoa} is usually traced to Latin \emph{macula} \textquotedblleft{}a spot, a stain\textquotedblright{} (Lewis \& Short), the same source as \emph{magoar}.
\end{cousinweb}

\begin{grammarlens}[title={What you've built}]
One more word for telling, explaining and describing at B1.
\end{grammarlens}

\subsection*{Your turn}
\begin{itemize}
\raggedright
  \item \emph{Say it:} \emph{a mágoa}
  \item \emph{Say it:} \emph{a mágoa}, once more
  \item \emph{Say it:} a sentence of your own with \emph{a mágoa}
  \item \emph{Recall:} say the Portuguese for grief, then the Portuguese for to hurt, then say \emph{a mágoa} again
\end{itemize}

\begin{tcolorbox}[breakable,colback=teal!4,colframe=teal!35!black,title={Before you move on}]
What does a mágoa mean? (\textquotedblleft{}Hurt, sorrow\textquotedblright{}.) Say it once more.
\end{tcolorbox}
\section[o rancor]{o rancor — grudge, resentment}

\label{lesson:PT-C479-rancor}



\begin{quote}
Before the new one: say the Portuguese for to hurt, then the Portuguese for hurt.
\end{quote}

\subsection*{You'll want to know: o rancor}
\textbf{o rancor} — \textquotedblleft{}grudge, resentment\textquotedblright{}.

\begin{quote}
Não guardo rancor a ninguém.
\end{quote}

\begin{quote}
\emph{I do not hold a grudge against anyone.}
\end{quote}

\emph{Guardar rancor a alguém} — to bear someone a grudge.

\begin{cousinweb}
\emph{rancor} comes from Late Latin \emph{rancor} \textquotedblleft{}rancidness, a grudge\textquotedblright{} (EtymDB), built on \emph{rancēre} \textquotedblleft{}to be rotten\textquotedblright{}.
\end{cousinweb}

\begin{grammarlens}[title={What you've built}]
One more word for telling, explaining and describing at B1.
\end{grammarlens}

\subsection*{Your turn}
\begin{itemize}
\raggedright
  \item \emph{Say it:} \emph{o rancor}
  \item \emph{Say it:} \emph{o rancor}, once more
  \item \emph{Say it:} a sentence of your own with \emph{o rancor}
  \item \emph{Recall:} say the Portuguese for to hurt, then the Portuguese for hurt, then say \emph{o rancor} again
\end{itemize}

\begin{tcolorbox}[breakable,colback=teal!4,colframe=teal!35!black,title={Before you move on}]
What does o rancor mean? (\textquotedblleft{}Grudge, resentment\textquotedblright{}.) Say it once more.
\end{tcolorbox}
\section[perdoar]{perdoar — to forgive}

\label{lesson:PT-C479-perdoar}



\begin{quote}
Before the new one: say the Portuguese for hurt, then the Portuguese for grudge.
\end{quote}

\subsection*{You'll want to know: perdoar}
\textbf{perdoar} — \textquotedblleft{}to forgive\textquotedblright{}.

\begin{quote}
A minha irmã perdoou-me o atraso.
\end{quote}

\begin{quote}
\emph{My sister forgave me for being late.}
\end{quote}

\emph{Perdoar alguma coisa a alguém} — to forgive someone for something.

\begin{cousinweb}
\emph{perdoar} is inherited from Old Portuguese \emph{perdõar} \textquotedblleft{}to forgive\textquotedblright{} (EtymDB), which is usually traced to Late Latin \emph{perdōnāre}, built on \emph{dōnāre} \textquotedblleft{}to give\textquotedblright{}.
\end{cousinweb}

\begin{grammarlens}[title={What you've built}]
One more word for telling, explaining and describing at B1.
\end{grammarlens}

\subsection*{Your turn}
\begin{itemize}
\raggedright
  \item \emph{Say it:} \emph{perdoar}
  \item \emph{Say it:} \emph{perdoar}, once more
  \item \emph{Say it:} a sentence of your own with \emph{perdoar}
  \item \emph{Recall:} say the Portuguese for hurt, then the Portuguese for grudge, then say \emph{perdoar} again
\end{itemize}

\begin{tcolorbox}[breakable,colback=teal!4,colframe=teal!35!black,title={Before you move on}]
What does perdoar mean? (\textquotedblleft{}To forgive\textquotedblright{}.) Say it once more.
\end{tcolorbox}
\section[o perdão]{o perdão — forgiveness, pardon}

\label{lesson:PT-C479-perdao}



\begin{quote}
Before the new one: say the Portuguese for grudge, then the Portuguese for to forgive.
\end{quote}

\subsection*{You'll want to know: o perdão}
\textbf{o perdão} — \textquotedblleft{}forgiveness, pardon\textquotedblright{}.

\begin{quote}
O rapaz pediu perdão à mãe.
\end{quote}

\begin{quote}
\emph{The boy asked his mother's forgiveness.}
\end{quote}

\emph{Pedir perdão} — to ask for forgiveness. \emph{Perdão?} — sorry, could you repeat that.

\begin{cousinweb}
\emph{perdão} comes from \emph{perdoar} \textquotedblleft{}to forgive\textquotedblright{} (EtymDB), from Late Latin \emph{perdōnāre}, built on \emph{dōnāre} \textquotedblleft{}to give\textquotedblright{}.
\end{cousinweb}

\begin{grammarlens}[title={What you've built}]
One more word for telling, explaining and describing at B1.
\end{grammarlens}

\subsection*{Your turn}
\begin{itemize}
\raggedright
  \item \emph{Say it:} \emph{o perdão}
  \item \emph{Say it:} \emph{o perdão}, once more
  \item \emph{Say it:} a sentence of your own with \emph{o perdão}
  \item \emph{Recall:} say the Portuguese for grudge, then the Portuguese for to forgive, then say \emph{o perdão} again
\end{itemize}

\begin{tcolorbox}[breakable,colback=teal!4,colframe=teal!35!black,title={Before you move on}]
What does o perdão mean? (\textquotedblleft{}Forgiveness, pardon\textquotedblright{}.) Say it once more.
\end{tcolorbox}
